\input{../shared/preamble}

\newcommand{\dissertationurl}{https://github.com/cderici/dissertation/releases/download/latest-pdf/thesis.pdf}
\newcommand{\pycketperfurl}{https://github.com/cderici/pycket-performance}
\newcommand{\tracedrawurl}{https://github.com/cderici/trace-draw}
\newcommand{\raxurl}{https://github.com/cderici/rax}
\newcommand{\icfpreporturl}{https://dl.acm.org/doi/10.1145/3341642}
\newcommand{\atwurl}{https://github.com/cderici/around-the-world-in-26-languages}
\newcommand{\homelaburl}{https://home.dericilab.live/homelab/}

\begin{document}

\paragraph{} %
Thank you for considering my application for compiler engineering roles at NVIDIA.

\paragraph{} %
My PhD work focused on building a self-hosting compiler on top of a meta-tracing JIT backend. That work required low-level performance engineering at the trace code generation level to address issues such as large traces creating instruction-cache pollution and GC pressure in the heap from tracing data-driven control flow. As part of this, I developed optimization techniques to reduce trace size, improved memory behavior by making key heap objects fit within L1 cache, and built performance analysis tools [\href{\dissertationurl}{1}, \href{\pycketperfurl}{2}, \href{\tracedrawurl}{3}]. I also hand-built a full Racket to x86\_64 assembly compiler with all the passes, including instruction selection, liveness analysis for register allocation via interference graphs, manual code generation for calling conventions and stack layout, and garbage collection [\href{\raxurl}{4}]. I also co-designed the intermediate representation currently shipped with the main Racket distribution, whose design and formal semantics are documented in my dissertation [\href{\icfpreporturl}{5}, \href{\dissertationurl}{1}].

\paragraph{} %
Since finishing my PhD, I have been developing my Around the World in 26 Languages project, where I experiment with a language/compiler for each letter of the alphabet [\href{\atwurl}{6}]. All the languages are written in C++, and most of them target LLVM IR, while some use MLIR dialects such as \emph{affine} and \emph{linalg} for optimization work based on polyhedral analysis. In one compilation path, I use the \texttt{nvptx64-nvidia-cuda} backend to generate PTX kernels targeting NVIDIA GPUs, and I use the CUDA Driver API to launch the generated kernels. More recently, I have been experimenting there with SIMT code generation, instruction pipelining, and loop parallelism through Tapir/LLVM and Cilk. Some of this work lives in my private homelab infrastructure, where I run experiments on actual NVIDIA GPUs on VMs [\href{\homelaburl}{7}]. In this project, I also experiment with autotuners to improve compilation performance using my machine learning background.

\paragraph{} %
During my PhD, I also worked at Canonical for three years, where I designed and developed a domain-specific language in Go for CPU and memory-constrained computation within the Juju cloud orchestration ecosystem. Aside from compilers, this gave me experience in distributed software, production software engineering, maintainability, collaboration, roadmap planning, hiring, and mentoring junior engineers.

\paragraph{} %
For these reasons, I believe my background is a strong fit for compiler engineering roles at NVIDIA. I’d be excited for the opportunity to contribute. Thank you again for your consideration.

\vspace{1.2em}

Sincerely,

Caner Derici

\vfill

{\small
\noindent\textbf{References} \\[0.25em]
\noindent \href{\dissertationurl}{[1] PhD Dissertation} \\
\noindent \href{\pycketperfurl}{[2] GitHub: pycket-performance} \\
\noindent \href{\tracedrawurl}{[3] GitHub: trace-draw} \\
\noindent \href{\raxurl}{[4] GitHub: Rax, Racket to x86\_64 Assembly} \\
\noindent \href{\icfpreporturl}{[5] Flatt, Derici, Dybvig, Keep et al., ``Rebuilding Racket on Chez Scheme'', ICFP 2019} \\
\noindent \href{\atwurl}{[6] GitHub: Around the World in 26 Languages} \\
\noindent \href{\homelaburl}{[7] Live Homelab Overview}
}

\vfill \hfill \small Last compiled on \today

\end{document}